\section{A minimal ambient-target Stokes shear}
\label{sec:incomplete-gamma}

Set \(\alpha=1/6\) and consider the rank-two meromorphic connection
\begin{equation}\label{eq:incomplete-gamma-connection}
 \frac{d}{dz}\binom{y_0}{y_1}=
 \begin{pmatrix}
 0&z^{-2}\\
 0&z^{-2}+\alpha z^{-1}
 \end{pmatrix}\binom{y_0}{y_1}.
\end{equation}

\begin{proposition}[Incomplete-Gamma countermodel]
\label{prop:incomplete-gamma}
Connection \eqref{eq:incomplete-gamma-connection} has a formal line of
monodromy \(\zetaSix\) whose sectorial
lift undergoes a nonzero Stokes shear into a monodromy-one line.  The shear
can be equipped with a flat hyperbolic pairing preserving a
\(\mathbf Z[\zetaSix]\)-lattice and can be tensored
with \(\Q[N]/(N^3)\) without disappearing.  Hence formal monodromy, pairing,
integrality, and a commuting length-three nilpotent tag do not determine the
zero or nonzero restriction of a flat covector to the \(\zetaSix\) line.
\end{proposition}

\begin{proof}
Two flat columns are
\begin{equation}\label{eq:incomplete-gamma-flat-columns}
 f_0=\binom10,
 \qquad
 f_\alpha(z)=\binom{\Gamma(1-\alpha,1/z)}{z^\alpha e^{-1/z}}.
\end{equation}
Indeed,
\begin{equation}\label{eq:incomplete-gamma-derivative}
 \frac d{dz}\Gamma(1-\alpha,1/z)=z^{\alpha-2}e^{-1/z},
\end{equation}
and direct differentiation verifies both rows of
\eqref{eq:incomplete-gamma-connection}.  The formal factors
are the trivial line and \(e^{-1/z}z^\alpha\), whose formal monodromy is
\(e^{2\pi i\alpha}=\zetaSix\).

Analytic continuation of the incomplete Gamma function adds a nonzero
multiple of the complete Gamma value.  After rescaling \(f_\alpha\), the
Stokes jump has the form
\begin{equation}\label{eq:incomplete-gamma-stokes-jump}
 f_\alpha^+=f_\alpha^-+f_0.
\end{equation}
Choose a flat covector \(r\) with \(r(f_0)=1\) and
\(r(f_\alpha^-)=0\).  Then \(r(f_\alpha^+)=1\).  Thus the Boolean restriction
of \(r\) to the sectorial lift of the same formal line changes.

Let \(A(z)\) be the matrix in \eqref{eq:incomplete-gamma-connection} and
double the system by its dual:
\begin{equation}\label{eq:hyperbolic-doubling}
 B(z)=\operatorname{diag}(A(z),-A(z)^t),\qquad
 J=\begin{pmatrix}0&I_2\\I_2&0\end{pmatrix}.
\end{equation}
Then \(B^tJ+JB=0\).  If the jump on the first summand is
\(S=\left(\begin{smallmatrix}1&1\\0&1\end{smallmatrix}\right)\),
the doubled jump \(\operatorname{diag}(S,S^{-t})\) preserves \(J\) and the
obvious \(\mathbf Z[\zetaSix]\)-lattice.  Finally tensor with
\(\Q[N]/(N^3)\) and let the jump act trivially on this factor.  It commutes
with \(N\) and still has \eqref{eq:incomplete-gamma-stokes-jump}.  This proves
the assertion.
\end{proof}

The example also occurs in a flat confluence family.  Replacing \(z\) by
\(z/t\) gives the \(z\)-equation
\begin{equation}\label{eq:incomplete-gamma-confluence}
 \frac{dY}{dz}=
 \begin{pmatrix}0&t z^{-2}\\0&t z^{-2}+\alpha z^{-1}\end{pmatrix}Y,
\end{equation}
with compatible \(t\)-equation inherited from the pullback connection.  The
two exponential factors collide at \(t=0\), while the nonzero Stokes
coefficient persists on the ramified punctured \(t\)-line.  Formal information
at the confluent fibre cannot recover its target orientation.

\begin{remark}
Proposition~\ref{prop:incomplete-gamma} is an analytic no-go theorem, not a
geometric counterexample.  We do not assert that
\eqref{eq:incomplete-gamma-connection} is the quantum
connection of a smooth projective wall crossing.  It proves only that a
geometric proof must use information absent from the listed formal shadows.
\end{remark}
